无创产前检测（NIPT）是一种通过采集母体血液检测胎儿游离 DNA 片段分析胎儿染色体是否异常从而达到产前诊断的技术方法，其核心判断依据为胎儿性染色体浓度：若孕期 10-25 周内，男胎 Y 染色体浓度达 4\% 及以上、女胎 X 染色体浓度无异常则NIPT结果基本准确，否则认为结果不具有准确性，需要进一步检测；胎儿异常发现时间与风险相关，12 周以内发现风险较低，13-27 周风险较高，28 周以后风险极高。男胎 Y 染色体浓度与孕妇孕周数、身体质量指数（BMI）密切相关$^{[7]}$ 。基于上述信息以及附件所包含的某地区孕妇的 NIPT 数据，本文需要解决以下问题：

\textbf{问题一}$\;$深入分析男胎孕妇中胎儿 Y 染色体浓度与孕周数、BMI 等指标之间的相关特性，构建能够量化所述变量关联的关系模型，并通过统计学方法检验该模型的显著性，以验证模型对变量间关系刻画的可靠性与有效性。

\textbf{问题二}$\;$男胎孕妇的 BMI 是影响胎儿 Y 染色体浓度达标时间的主要因素。对男胎孕妇的 BMI 进行合理分组，明确每组具体的 BMI 区间，同时结合尽可能早地发现胎儿异常的目标，确定每组对应的最佳 NIPT 时点，并分析检测误差对分组结果的影响。

\textbf{问题三}$\;$男胎 Y 染色体浓度达标时间不仅受 BMI 影响，还与孕妇身高、体重、年龄等多种因素相关。现综合考虑这些因素，结合检测误差以及胎儿 Y 染色体浓度达标比例，对男胎孕妇的 BMI 进行合理分组，确定每组的最佳 NIPT 时点。最后同样需分析检测误差对分组合理性的影响。

\textbf{问题四}$\;$女胎孕妇的 21 号、18 号、13 号染色体非整倍体判定结果为核心依据，综合考虑 X 染色体及上述三种染色体的 Z 值、GC 含量、读段数及相关比例、孕妇 BMI 等因素，构建女胎异常判定方法。 
